\documentclass[11pt]{article}
\usepackage[a4paper,margin=0.92in]{geometry}
\usepackage{amsmath,amssymb,amsthm,booktabs,microtype}
\usepackage[hidelinks]{hyperref}
\newtheorem{proposition}{Proposition}
\title{Two-Way Log-Budget Decomposition of\
Prime-Shell Operator Interaction}
\author{Liang Wang\\
School of Mathematics and Statistics\\
Huazhong University of Science and Technology (HUST), Wuhan, China}
\date{29 August 2026}

\begin{document}
\maketitle

\begin{abstract}
TPC-305 replaced one shell's source-first target by an aligned neighboring
label while holding that shell's physical operator fixed.  The resulting
orientation atlas isolated a within-row target comparison but left the
two-row operator interaction implicit.  We arrange the four positive budget
cells as an operator-by-target table and define two log target-switch effects.
Their mean is a target-main contrast $m$, their half-difference is an operator
interaction contrast $i$, and the exact identity
$m^2-i^2=d_Ld_R$ gives a strict dominance test.  Replaying all 18 TPC-305
cases at three source normalizers yields 54 derived rows: target-main
dominance occurs in 12 of 18 cases and interaction dominance in 6 of 18, with
no unresolved case.  At the central $Q=60\to70$ fracture the split is 5/6
versus 1/6, and all three inherited same-prefix cases are target-main
dominant.  The finite diagnostic quantifies, but does not remove, the
common-ambient and causal-identification obstruction.
\end{abstract}

\section{From a target swap to a two-way table}

TPC-305 is the immediate predecessor of this report \cite{tpc305}.  It holds each physical
shell/operator fixed, changes the target only on an overlap using the aligned
neighboring sign label, and retains the native sign off the overlap.  On the
fixed $(N,H,z)=(512,58,5)$ spine with $Q=50,60,70,90$, it finds five of six
central cases in which the right-neighbor target is cheaper on both fixed
operators.  That statement is useful, but two operators still have different
physical output rows.  A single orientation label therefore hides an
interaction between operator identity and target identity.

We retain exactly the TPC-305 finite data: two kernel exponents, three
relative RMS tolerances, three adjacent shell pairs, and three positive source
normalizers.  No new shell or asymptotic regime is introduced.  For one pair,
write the four normalized budget cells as
\[
\begin{array}{c|cc}
 & \text{left target} & \text{right target}\\ \hline
\text{left operator} & B_{LL} & B_{LR}\\
\text{right operator} & B_{RL} & B_{RR}
\end{array}                                                     \tag{1}
\]
Here ``left target'' means the native left label on the left row and its
aligned transported completion on the right row; ``right target'' has the
opposite completion.  Thus (1) is a protocol-level holdout table, with the
same shell-specific completion convention as TPC-305.

\section{Effects and exact decomposition}

The common intervention is the target switch left $\to$ right.  Its row-wise
log effects are
\[
 d_L=\log\frac{B_{LR}}{B_{LL}},\qquad
 d_R=\log\frac{B_{RR}}{B_{RL}}.                                  \tag{2}
\]
TPC-305 recorded $R_L=B_{LR}/B_{LL}$ and
$R_R=B_{RL}/B_{RR}$, so $d_L=\log R_L$ and $d_R=-\log R_R$.  Define
\[
 m=\frac{d_L+d_R}{2},\qquad
 i=\frac{d_L-d_R}{2},\qquad
 q=\frac{|i|}{|m|}.                                               \tag{3}
\]
The quantity $m$ measures the target-switch effect averaged over the two
operator rows; $i$ measures how much that effect changes between rows.

\begin{proposition}[Two-way squared-contrast identity]
For positive cells in (1),
\[
 d_L=m+i,\qquad d_R=m-i,\qquad m^2-i^2=d_Ld_R.                   \tag{4}
\]
If both row effects are nonzero, equal signs are equivalent to $|m|>|i|$,
and opposite signs are equivalent to $|i|>|m|$.
\end{proposition}
\begin{proof}
Adding and subtracting (3) gives the first two identities.  Expanding the
difference of squares gives
\[
 m^2-i^2=\frac{(d_L+d_R)^2-(d_L-d_R)^2}{4}=d_Ld_R.
\]
The product on the right is positive exactly for equal nonzero signs and
negative exactly for opposite signs.  Taking square roots gives the two strict
inequalities.
\end{proof}

\begin{proposition}[Row-normalization invariance]
Independent positive rescalings of the two rows of (1) leave $d_L,d_R,m,i$
and $q$ unchanged.
\end{proposition}
\begin{proof}
Multiplying both cells in one row by the same positive factor cancels in its
ratio in (2).  Apply this separately to the two rows.
\end{proof}

We call $|m|>|i|$ target-main dominance and $|i|>|m|$ interaction dominance.
The sign of the two effects distinguishes the right-neighbor and left-neighbor
main preferences.  This terminology is algebraic; it is not a causal claim.

\section{Finite certificate}

The producer reads the locked TPC-305 ratio intervals and applies 80-digit
logarithmic interval arithmetic with a padding of $10^{-30}$.  The independent
checker does not import this producer: it replays the logarithms with
100-digit Decimal arithmetic, verifies all 54 enclosures, checks (4) on the
printed centers, and reconstructs the case census.  A stress suite exhausts
625 positive four-cell tables and independent positive row scalings.

Table~\ref{tab:pair} gives the case-level result.  A main-dominant case has
same-sign row effects; an interaction-dominant case has opposite signs.

\begin{table}[t]
\centering\small
\caption{Two-way dominance census over six cases per transition.}
\label{tab:pair}
\begin{tabular}{c c c c}
\toprule
transition & target-main & interaction & same-prefix main\\
\midrule
$50\to60$ & 4 & 2 & 0\\
$60\to70$ & 5 & 1 & 3\\
$70\to90$ & 3 & 3 & 0\\
\bottomrule
\end{tabular}
\end{table}

The central six rows are displayed in Table~\ref{tab:central}.  The numbers
are interval centers; the certificate stores 38-digit endpoints.  The three
normalizers agree on every strict class.

\begin{table}[t]
\centering\small
\caption{Central log effects and contrasts (interval centers).}
\label{tab:central}
\begin{tabular}{c c r r r l}
\toprule
$e$ & $\tau$ & $m$ & $i$ & $q$ & dominant term\\
\midrule
1 & .25 & $-.203481$ & $.500980$ & $2.46205$ & interaction\\
1 & .50 & $-.209604$ & $-.133643$ & $.637596$ & target-main\\
1 & .75 & $-1.27333$ & $-.629188$ & $.494130$ & target-main\\
2 & .25 & $-1.22433$ & $-1.07622$ & $.879029$ & target-main\\
2 & .50 & $-2.20188$ & $-.636939$ & $.289271$ & target-main\\
2 & .75 & $-3.19790$ & $-1.85187$ & $.579089$ & target-main\\
\bottomrule
\end{tabular}
\end{table}

Across all 36 target-main normalizer rows, the upper endpoint of $q$ is below
$0.88$.  Across all 18 interaction-dominant normalizer rows, the lower
endpoint is above $1.2$.  For the three central same-prefix cases the maximum
upper endpoint is below $0.64$.  Thus the finite atlas has a visible gap
between the two algebraic regimes rather than a borderline classification.

\section{What this does and does not identify}

The positive result is a compact interaction diagnostic.  TPC-305's central
five-of-six right-label signal is exactly the five-of-six target-main regime
in the table, and the same-prefix parent descents remain target-main dominant.
The identity (4) explains why: both row-wise target effects have the same
sign, so their mean is larger in magnitude than their half-difference.

The obstruction is not erased.  Six cases have opposite-sign row effects, so
operator interaction is larger than the mean target contrast; one of these is
at the central fracture and three occur on the final transition.  Moreover,
the right and left targets in (1) use native off-overlap completions specific
to their row.  Consequently the table is not a balanced intervention on a
single common ambient target vector.  A common union-shell completion and an
operator holdout are required before interpreting $m$ as a causal effect.

The decomposition also supplies no arithmetic $L^2$ estimate, fixed power
saving, full Route-B Gate-B theorem, or twin-prime result.  The Session-named
Route-A/Route-B evaluator files are absent from this checkout, so the claim
level is the local fail-closed level recorded in the route-evaluation note.

\section*{Reproducibility and next question}

The canonical certificate is
\texttt{results/tpc306\_certificate.json}.  The producer, independent checker,
stress suite, and Bridge-B checker are included with the project.  The next
minimal question is whether the dominance gap survives a common-ambient
union-shell completion and alternative off-overlap transports.  No claim is
made that the finite orientation persists as $Q$ or $N$ grows.

\bibliographystyle{plain}
\bibliography{references}
\end{document}
